% Converted from reference/DB4OFlare (MadCap Flare, Java target)

\chapter{Community}
\label{ref:community}

Participate in the db4o community and help to improve db4o. The main community website for db4o is \href{http://community.versant.com/}{http://community.versant.com/}. There you can find the db4o forums, blog posts additional resources and the newest db4o releases.

\section{Forums}
\label{ref:community:db4o}

The best place to ask questions about db4o, discuss and make suggestions are the db4o forums. You can find the db4o forums here:\href{http://community.versant.com/Forums.aspx}{http://community.versant.com/Forums.aspx}

\section{Resources \& News}
\label{ref:community:addition}

To stay informed about new features, community activity and other news subscribe to the different db4o blogs: \href{http://community.versant.com/Blogs.aspx}{http://community.versant.com/Blogs.aspx}

Or follow us on Twitter: \href{http://twitter.com/db4objects}{http://twitter.com/db4objects}

\section*{Report Bugs}

Every software contains bugs and so does db4o. If you find bugs please report them to us. In case you're not sure if you really found a bug, you maybe want to discuss your issue first on the \href{http://community.versant.com/Forums.aspx}{db4o forum}. As soon as you're sure that you found a bug it best to report in \href{http://tracker.db4o.com/}{our bug-tracking system}. The login is the same as in the db4o forums. Read more on tips and tricks for reporting bugs \hyperref[ref:community-report-bugs]{here} (p.~\pageref{ref:community-report-bugs}).

When possible include a small test program which reproduces the issue with the bug report. A test case avoids confusion and effort to reproduce the bug.

\section*{Working with the Source Code}

\hyperref[ref:community-working-with-source-code]{See "Working With Source Code"} (p.~\pageref{ref:community-working-with-source-code})

\section{Report Bugs}
\label{ref:community-report-bugs}

Found a bug in db4o? Here are a few tips and tricks for reporting bugs.

\subsection*{Where Do I Report Bugs?}

The best way to report a bug is to report it on our issue \href{http://tracker.db4o.com/}{tracking system}. The login for the bug tracker is the same as in the db4o community page and db4o forums. If you don't have a login yet you can register \href{http://community.versant.com/register.aspx}{here}.

In cases you are not sure if the issue is an bug or another problem you maybe want to discuss it in the \href{http://community.versant.com/register.aspx}{db4o forums}.

\subsection*{What Should the Bug Report Contain?}

The bug report should be as specific as possible, so that the room for interpretation is very small:

\begin{itemize}
\item Steps to reproduce the bug. In order to fix the bug a step by step instruction on how to reproduce the issue is essential.
\item When possible include a the test-code which reproduces the issue.
\item Include the db4o version in your bug report. The complete version of db4o is included in the jar-name.
\item Include information about your platform. Which Java Version and what kind of device (like Android).  Which .NET version and what kind of device (for example a embedded system)
\item In case of exceptions, please add also the stack trace.
\end{itemize}

\section{Working With Source Code}
\label{ref:community-working-with-source-code}

db4o is an open-source project. The source is available for reviewing, modifying for own needs or contributing your modifications. You can use the \hyperref[ref:community-working-with-source-using-the-sources-from-db4o-distribution]{source code from the downloaded distribution} (p.~\pageref{ref:community-working-with-source-using-the-sources-from-db4o-distribution}) package or you may use our \hyperref[ref:community-working-with-source-using-the-repository]{SVN repository} (p.~\pageref{ref:community-working-with-source-using-the-repository}) to get the latest modifications.

More Reading:

\begin{itemize}
\item \hyperref[ref:community-working-with-source-db4o-testing-framework]{Unit Tests for db4o} (p.~\pageref{ref:community-working-with-source-db4o-testing-framework})
\item \hyperref[ref:community-working-with-source-sharpen]{Sharpen: Converting Java to C\#} (p.~\pageref{ref:community-working-with-source-sharpen})
\item \hyperref[ref:community-working-with-source-using-the-sources-from-db4o-distribution]{Using The Sources From db4o Distribution} (p.~\pageref{ref:community-working-with-source-using-the-sources-from-db4o-distribution})
\item \hyperref[ref:community-working-with-source-using-the-repository]{Using The Repository} (p.~\pageref{ref:community-working-with-source-using-the-repository})
\item \hyperref[ref:community-working-with-source-building-full-distribution]{Building Full Distribution} (p.~\pageref{ref:community-working-with-source-building-full-distribution})
\item Building Java Version
\item \hyperref[ref:community-working-with-source-patch-submission]{Patch Submission} (p.~\pageref{ref:community-working-with-source-patch-submission})
\end{itemize}

\subsection{Sharpen}
\label{ref:community-working-with-source-sharpen}

Sharpen is an tool that allows you to convert your Java db4o project into C\#. The difference between sharpen and other java-to-C\# converters is the native support for db4o, .NET naming conventions and customization options.

More Reading:

\begin{itemize}
\item \hyperref[ref:community-working-with-source-sharpen-how-to-setup-sharpen]{How To Setup Sharpen} (p.~\pageref{ref:community-working-with-source-sharpen-how-to-setup-sharpen})
\item \hyperref[ref:community-working-with-source-sharpen-example-conversion]{Doing a Sharpen Conversion} (p.~\pageref{ref:community-working-with-source-sharpen-example-conversion})
\item \hyperref[ref:community-working-with-source-sharpen-sharpen-command-line-arguments]{Sharpen Command-Line Arguments} (p.~\pageref{ref:community-working-with-source-sharpen-sharpen-command-line-arguments})
\item \hyperref[ref:community-working-with-source-sharpen-sharpen-annotations]{Sharpen Annotations} (p.~\pageref{ref:community-working-with-source-sharpen-sharpen-annotations})
\end{itemize}

\subsubsection{How To Setup Sharpen}
\label{ref:community-working-with-source-sharpen-how-to-setup-sharpen}

You can obtain sharpen source from db4o svn repository at:

\href{https://source.db4o.com/db4o/trunk/sharpen}{https://source.db4o.com/db4o/trunk/sharpen}

For the ease of use check-out sharpen projects:

\begin{itemize}
\item sharpen.builder
\item sharpen.core
\item sharpen.ui
\item sharpen.ui.tests
\end{itemize}

Additionally Sharpen requires a valid \href{http://eclipse.org/}{Eclipse} installation to run. Install Eclipse on your machine or reuse an existing installation. Check out the four projects into the same workspace as your java project which you want to convert. This is not required is easier to use and maintain.

The first step is to build sharpen. For that we can use Apache Ant:

\begin{lstlisting}
<target name="build-sharpen">
    <property name="sharpen.core.dir" location="${dir.workspace}/sharpen.core"/>
    <reset-dir dir="${dir.dist.classes.sharp}"/>

    <javac fork="true"
           debug="true"
           target="1.5"
           source="1.5"
           destdir="${dir.dist.classes.sharp}"
           srcdir="${sharpen.core.dir}/src"
           encoding="UTF-8">
        <classpath>
            <fileset dir="${eclipse.home}/plugins">
                <include name="org.eclipse.osgi_*.jar"/>
                <include name="org.eclipse.core.resources_*.jar"/>
                <include name="org.eclipse.core.runtime_*.jar"/>
                <include name="org.eclipse.jdt.core_*.jar"/>
                <include name="org.eclipse.jdt.launching_*.jar"/>
                <include name="org.eclipse.equinox.*.jar"/>
                <include name="org.eclipse.core.jobs_*.jar"/>
            </fileset>
        </classpath>
    </javac>
    <jar destfile="${dir.dist.classes.sharp}/sharpen.core_1.0.0.jar" basedir="${dir.dist.classes.sharp}">
        <fileset dir="${sharpen.core.dir}">
            <include name="plugin.xml"/>
        </fileset>
    </jar>
</target>
\end{lstlisting}

{\small\itshape sharpen-install.xml: Build Sharpen\par}
\medskip

To run Sharpen you should install it to an Eclipse instance:

\begin{lstlisting}
<target name="install-sharpen-plugin" depends="build-sharpen">
    <copyfile src="${dir.dist.classes.sharp}/sharpen.core_1.0.0.jar" dest="${plugins.home}/sharpen.core_1.0.0.jar"/>
</target>
\end{lstlisting}

{\small\itshape sharpen-install.xml: Install Sharpen to Eclipse\par}
\medskip

Put the paths for the build in a property file, so that you can easily change them. Here's a example of the property file. You have to configure the JDK-path, the Eclipse path and the path to the Sharpen source.

\begin{lstlisting}
#The workspace where the sharpen projects are
dir.workspace=C:/Users/Gamlor/Develop/db4o/sharpenProject
# Java executable
jdk.home=${env.JAVA_HOME}
jdk.home.java=${jdk.home}/bin/java.exe
# Eclipse home directory
eclipse.home=C:/progs/eclipse
# Sandcastle can be used to convert javadoc to .NET xml comments
# dir.lib.sandcastle=e:/sandcastle/
# sharpen compile directory
dir.dist.classes.sharp=${dir.workspace}/dist/
# Eclipse plugins home
plugins.home=${eclipse.home}/plugins
\end{lstlisting}

{\small\itshape sharpen.properties: The configuration for building sharpen\par}
\medskip

After that Sharpen is set up to run.

\subsubsection{Doing a Sharpen Conversion}
\label{ref:community-working-with-source-sharpen-example-conversion}

Ensure that you've installed sharpen to an existing eclipse installation as explained \hyperref[ref:community-working-with-source-sharpen-how-to-setup-sharpen]{here} (p.~\pageref{ref:community-working-with-source-sharpen-how-to-setup-sharpen}).

Use Ant scripts to run Sharpen and translate your Java code to C\#. The best way for this is to define an Ant macro which you then can reuse. This task takes two arguments. The first argument is the path to a valid Eclipse workspace which contains the project to translate. The second parameter is the project in the workspace which you want to translate.

\begin{lstlisting}
<macrodef name="sharpen">
    <attribute name="workspace"/>
    <attribute name="resource"/>

    <element name="args" optional="yes"/>

    <sequential>
        <java taskname="sharpen"
              fork="true"
              classname="org.eclipse.core.launcher.Main"
              failonerror="true" timeout="1800000">

            <classpath>
                <fileset dir="${eclipse.home}/plugins">
                    <include name="org.eclipse.equinox.launcher_*.jar"/>
                </fileset>
            </classpath>

            <arg value="-clean"/>
            <arg value="-data"/>
            <arg file="@{workspace}"/>
            <arg value="-application"/>
            <arg value="sharpen.core.application"/>
            <arg value="@{resource}"/>
            <args/>
        </java>
    </sequential>
</macrodef>
\end{lstlisting}

{\small\itshape sharpen-install.xml: The sharpen task\par}
\medskip

Now you can use this task to sharpen your project. First ensure that your project is in a valid Eclipse workspace. Then you specify the workspace and the sources of the project:

\begin{lstlisting}
<target name="sharpen">
    <sharpen
        workspace="C:\temp\sharpenExamples\"
        resource="example/src">
        <args>
            <arg value="@sharpen-config"/>
        </args>
        </sharpen>
</target>
\end{lstlisting}

{\small\itshape sharpen-example.xml: Sharpen a example project\par}
\medskip

Additionally you can pass the sharpen configuration as a file-name. When you add a '@' in front of the file-name sharpen will read that file and use all configuration flags of that. For example:

\begin{lstlisting}
-pascalCase+
-nativeTypeSystem
-nativeInterfaces
\end{lstlisting}

You can find a list of all Sharpen \hyperref[ref:community-working-with-source-sharpen-sharpen-command-line-arguments]{configuration flags here} (p.~\pageref{ref:community-working-with-source-sharpen-sharpen-command-line-arguments}) and a list of all \hyperref[ref:community-working-with-source-sharpen-sharpen-annotations]{Sharpen annotations here} (p.~\pageref{ref:community-working-with-source-sharpen-sharpen-annotations}).

\subsubsection{Sharpen Command-Line Arguments}
\label{ref:community-working-with-source-sharpen-sharpen-command-line-arguments}

Sharpen command-line arguments can be defined in an options file

\begin{lstlisting}
<sharpen workspace="${target.dir}" resource="sharpened_examples/src">
   <args>
      <!-- Sharpen options are defined in a separate file -->
      <arg value="@sharpen-all-options" />
    </args>
 </sharpen>      
\end{lstlisting}

Here sharpen-all-options file contains all command-line options needed to convert current project. For an example of command-line options file see the \hyperref[ref:community-working-with-source-sharpen-example-conversion]{previous topic} (p.~\pageref{ref:community-working-with-source-sharpen-example-conversion}).

Command-line arguments can also be specified directly in an ant script:

\begin{lstlisting}
<sharpen workspace="${dir.sharpen}" resource="db4oj/core/src">
    <args>
        <arg value="-xmldoc"/>
        <arg file="config/sharpen/ApiOverlay.xml" />       
        <arg value="@sharpen-all-options" />
     </args>
 </sharpen>
        
\end{lstlisting}

The following table shows available command-line options, their meaning and example usage:

{\small
\begin{longtable}{|p{\dimexpr(\linewidth-3\arrayrulewidth)*1/2-2\tabcolsep\relax}|p{\dimexpr(\linewidth-3\arrayrulewidth)*1/2-2\tabcolsep\relax}|}
\hline
\textbf{Argument} & \textbf{Usage} \\ \hline
-pascalCase & Convert Java identifiers to Pascal case \\ \hline
-pascalCase+ & Convert Java indentifiers and package names (namespaces) to Pascal case \\ \hline
-cp & Adds a new entry to classpath:\par \textless{}arg value="-cp" /\textgreater{}\par \textless{}arg path="lib/db4o-7.2.37.10417-java5.jar" /\textgreater{} \\ \hline
-srcFolder & Adds a new source folder for sharpening \\ \hline
-nativeTypeSystem & Map java classes to .NET classes with a similar functionality. For example: java.lang.Class - System.Type \\ \hline
-nativeInterfaces & Adds an "I" in front of the interface name \\ \hline
-organizeUsings & Adds "using" for the types used \\ \hline
-fullyQualify & Converts to a fully-qualified name:\par -fullyQualify File \\ \hline
-namespaceMapping & Maps a java package name to a .NET namespace. For example:\par -namespaceMapping com.db4o Db4objects.Db4o \\ \hline
-methodMapping & Maps a java method name to a .NET method (can be method in another class). For example:\par -methodMapping java.util.Date.getTime Sharpen.Runtime.ToJavaMilliseconds \\ \hline
-typeMapping & Maps a java class to .NET type:\par -typeMapping com.db4o.Db4o Db4objects.Db4o.Db4oFactory \\ \hline
-propertyMapping & Maps a java method to .NET property:\par -propertyMapping com.db4odoc.structured.Car.getPilot Pilot \\ \hline
-runtimeTypeName & Name of the runtime class. The runtime class provides implementation for methods that don't have a direct mapping or that are simpler to map at the language level than at the sharpen level. For instance: String.substring, String.valueOf, Exception.printStackTrace, etc.\par For a complete list of all the method that can be mapped to the runtime class see Configuration\#runtimeMethod call hierarchy. \\ \hline
-header & Header comment to be added to all converted files.\par -header config/copyright\_comment.txt \\ \hline
-xmldoc & Specifies an xml-overlay file, which overrides javadoc documentation for specific classes:\par -xmldoc config/sharpen/ApiOverlay.xml \\ \hline
-eventMapping & Converts the methods to an event. \\ \hline
-eventAddMapping & Marks the method as an event subscription method. Invocations to the method in the form \textless{}target\textgreater{}.method(\textless{}argument\textgreater{}) will be replaced by the c\# event subscription idiom: \textless{}target\textgreater{} += \textless{}argument\textgreater{} \\ \hline
-conditionalCompilation & Add a condition when to translate the Java code\par - -conditionalCompilation com.db4o.db4ounit.common.cs !SILVERLIGHT \\ \hline
-configurationClass & Change the configuration class. The default is 'sharpen.core.DefaultConfiguration' \\ \hline
\end{longtable}}

\subsubsection{Sharpen Annotations}
\label{ref:community-working-with-source-sharpen-sharpen-annotations}

Sharpen annotations decorate java source code and are used to notify sharpener about how the code should be processed and converted. Annotations can be used to specify how a code element should be converted (for example class to enum), to skip conversion of some code elements, to rename classes, to change visibility etc.

The following table shows existing annotations, their meaning and examples.

{\small
\begin{longtable}{|p{\dimexpr(\linewidth-3\arrayrulewidth)*1/2-2\tabcolsep\relax}|p{\dimexpr(\linewidth-3\arrayrulewidth)*1/2-2\tabcolsep\relax}|}
\hline
\textbf{Annotation} & \textbf{Meaning} \\ \hline
@sharpen.enum & Mark java class to be processed as a .NET enum \\ \hline
@sharpen.rename & Specifies a different name for the converted type, takes a single name argument. For example:\par @sharpen.rename Db4oFactory \\ \hline
@sharpen.private & Specifies that the element must be declared private in the converted file, though it can be not private in the java source:\par \begin{lstlisting}
/*
* @sharpen.private
*/
\end{lstlisting}\par public List4 \_first; \\ \hline
@sharpen.internal & Specifies that the element must be declared internal in the converted file:\par \begin{lstlisting}
/**
 * @sharpen.internal
*/
\end{lstlisting}\par public abstract int size(); \\ \hline
@sharpen.protected & Specifies that the element must be declared protected in the converted file:\par \begin{lstlisting}
/**
 * @sharpen.protected
*/
\end{lstlisting}\par public abstract int size(); \\ \hline
@sharpen.new & Adds the C\#-'new' modifier to the translated code. \\ \hline
@sharpen.event & Links an event to its arguments. For example:\par Java:\par \begin{lstlisting}
/**
* @sharpen.event com.db4o.events.QueryEventArgs
*/
public Event4 queryStarted();
\end{lstlisting}\par is converted to:\par \begin{lstlisting}
public delegate void QueryEventHandler(
  object sender,
  Db4objects.Db4o.Events.QueryEventArgs args);
.......
event Db4objects.Db4o.Events.QueryEventHandler QueryStarted;
\end{lstlisting} \\ \hline
@sharpen.event.add & Marks the method as an event subscription method. Invocations to the method in the form \textless{}target\textgreater{}.method(\textless{}argument\textgreater{}) will be replaced by the c\# event subscription idiom: \textless{}target\textgreater{} += \textless{}argument\textgreater{} \\ \hline
@sharpen.event.onAdd & Valid for event declaration only (SHARPEN\_EVENT). Configures the method to be invoked whenever a new event handler is subscribed to the event. \\ \hline
@sharpen.if & Add \#if \textless{}expression\textgreater{}\#endif declaration:\par @sharpen.if \textless{}expression\textgreater{} \\ \hline
@sharpen.property & Convert a java method as a property:\par \begin{lstlisting}
/**
 * @sharpen.property
*/
\end{lstlisting}\par public abstract int size(); \\ \hline
@sharpen.indexer & Marks an element as an indexer property \\ \hline
@sharpen.ignore & Skip the element while converting \\ \hline
@sharpen.ignore.extends & Ignore the extends clause in Java class definition \\ \hline
@sharpen.ignore.implements & Ignore the implements clause in Java class definition \\ \hline
@sharpen.extends & Adds an extends clause to the converted class definition. For example:\par Java:\par \begin{lstlisting}
/**
* @sharpen.extends System.Collections.IList
*/
public interface ObjectSet {...
\end{lstlisting}\par converts to\par \begin{lstlisting}
public interface IObjectSet : System.Collections.IList
\end{lstlisting} \\ \hline
@sharpen.partial & Marks the converted class as partial \\ \hline
@sharpen.remove & Marks a method invocation that should be removed \\ \hline
@sharpen.remove.first & Removes the first line of the method/constructor when converting to C\#:\par \begin{lstlisting}
/**
* @sharpen.remove.first
*/
public void doSomething(){
    System.out.println("Java");
    NextMethod();
}
\end{lstlisting}\par converts to:\par \begin{lstlisting}
public void DoSomething(){
    NextMethod();
}
\end{lstlisting} \\ \hline
@sharpen.struct & Marks class to be converted as c\# struct \\ \hline
@sharpen.unwrap & When a method is marked with this annotation \newline{}all method calls are removed. This is useful for removing conversion methods when their aren't required in C\#.\par \begin{lstlisting}
/*
* @sharpen.unwrap
*/
public Iterable toIterable(Object[] array){
   return Arrays.asList(array);
}
public void doSomething(Object[] objs){
  Iterable iterable = toIterable(objs);
  // do something with the iterable
}
\end{lstlisting}\par Is converted to:\par \begin{lstlisting}
public IEnumerable ToIterable(object[] array){
    return Arrays.AsList(array);
}
public void doSomething(object[] objs){
   Iterable iterable = objs;
   // do something with the iterable
}
\end{lstlisting} \\ \hline
@sharpen.attribute & Adds an attribute to the converted code:\par \begin{lstlisting}
/*
* @sharpen.attribute TheAttribute
*/
public void doSomething(){}
\end{lstlisting}\par Will be converted to:\par \begin{lstlisting}
[TheAttribute]
public void DoSomething(){}
\end{lstlisting} \\ \hline
@sharpen.macro & Add a replace-pattern macro to your code. \\ \hline
\end{longtable}}

\subsection{db4o Testing Framework}
\label{ref:community-working-with-source-db4o-testing-framework}

db4ounit is a minimal xUnit (jUnit, NUnit) style testing framework. The db4ounit framework was created to fulfill the following requirements:

\begin{itemize}
\item The core tests should be run against JDK1.1
\item It should be possible to automatically convert test cases from Java to .NET.
\end{itemize}

db4ounit design deviates from vanilla xUnit in some respect, but if you know xUnit, db4ounit should look very familiar.

db4ounit itself is completely agnostic of db4o, but there is the db4ounit.extensions module which provides a base class for db4o specific test cases with different fixtures, etc.

Db4ounit and db4ounit.extensions are supplied as a source code for both java and .NET. Java version also comes with a compiled library: db4o-X.XX-db4ounit.jar, which allows you to run your tests from a separate package.

If you've found a bug and want to supply a test case to help db4o to fix the issue quickly, the best option would be to supply your code in the java db4ounit format. This format allows very easy integration of a new test case into db4o test suite: only copy/paste is required to put your test class code into the framework using Eclipse.

More Reading:

\begin{itemize}
\item \hyperref[ref:community-working-with-source-db4o-testing-framework-creating-a-sample-test]{Creating A Sample Test} (p.~\pageref{ref:community-working-with-source-db4o-testing-framework-creating-a-sample-test})
\item \hyperref[ref:community-working-with-source-db4o-testing-framework-db4ounit-methods]{Db4ounit Methods} (p.~\pageref{ref:community-working-with-source-db4o-testing-framework-db4ounit-methods})
\end{itemize}

\subsubsection{Creating A Sample Test}
\label{ref:community-working-with-source-db4o-testing-framework-creating-a-sample-test}

Let's create the a extremely simple test case. The first step is to setup db4ounit. The best way is to use db4ounit directly from the source. db4ounit is in the db4o source code folder of the distribution.

Open the db4o source-code of the db4o test-projects in Eclipse (or another IDE). The db4ounit projects are called 'db4ounit' and 'db4ounit.extensions'. Then create a new project which references the db4ounit projects.

After that we're ready to write our first db4ounit test. Create a new class which inherits from AbstractDb4oTestCase. Then add a test-method. Any method which starts with the prefix 'test' is a test method. Then add a main method which starts the test.

\begin{lstlisting}
public class ExampleTestCase extends AbstractDb4oTestCase{

    public static void main(String[] args) {
        new ExampleTestCase().runEmbedded();
    }

    public void testStoresElement(){
        db().store(new TestItem());
        ObjectSet<TestItem> result = db().query(TestItem.class);
        Assert.areEqual(1, result.size());
    }

    static class TestItem{

    }
}
\end{lstlisting}

{\small\itshape ExampleTestCase.java: Basic test case\par}
\medskip

\subsubsection{db4ounit Methods}
\label{ref:community-working-with-source-db4o-testing-framework-db4ounit-methods}

Let's look through the basic API , which will help you to build your own test. This document is not a complete API reference and its intention is to give you a general idea of the methods usage and availability.

\paragraph*{AbstractDb4oTestCase}

AbstractDb4oTestCase is a base class for creating test cases. It will setup a db4o instance for you which you can use in your tests. Additionally it provides different utility methods for configuring, querying and modifying the database.

Additionally it contains methods to run the test. You can create new instance of a class which extends AbstractDb4oTestCase and run the test in different environments. For example the method 'runSolo' will run the test with an embedded local container.

\paragraph*{TestCase}

This interface provides the basic for a unit test in the db4ounit framework. When you implement this interface the class is a valid unit test. All methods starting with 'test' will be executed as test.

For most tests it is more continent to use the AbstractDb4oTestCase.

\paragraph*{TestLifeCycle}

This interface provides a test case with additional setup and tear down methods. Those will be called before and after each test method

\paragraph*{ConsoleTestRunner}

A test-runner which runs the tests as console application.

Usually it's more convenient to extend the AbstractDb4oTestCase class and use the provided run methods instead of the console test runner.

\paragraph*{Db4ounit.Assert}

This class provides a variety of methods for asserting certain conditions

\subsection{Using The Sources From db4o Distribution}
\label{ref:community-working-with-source-using-the-sources-from-db4o-distribution}

The sources for db4o are in the 'src' subdirectory of the db4o distribution. You simply can open these projects in Eclipse (or another IDE). Compile and use the db4o projects like any other Eclipse project.

Following projects are shipped with db4o:

\begin{itemize}
\item bloat: The bytecode manipulation library used by db4o.
\item db4o.cs: The db4o client server mode.
\item db4o.cs.optional: Optional additions for the db4o client server mode.
\item db4o.instrumentation: The base for different instrumentations provided by db4o.
\item db4o\_osgi: The db4o OSGi service
\item db4o\_osgi\_test: The test cases for the db4o OSGi service.
\item db4oj: The db4o core.
\item db4oj.optional: Optional features for db4o.
\item db4oj.tests: The db4o test-suite.
\item db4onqopt: The db4o native query optimization.
\item db4otaj: Transparent activation/persistence enhancers.
\item db4otools: Tooling for db4o like Ant-tasks
\item db4ounit: The db4o unit test framework.
\item db4ounit.extensions: Utilities and extensions for the db4o unit test framework.
\end{itemize}

\subsection{Using The Repository}
\label{ref:community-working-with-source-using-the-repository}

If you enjoy being on the "cutting edge" and want to follow up with the development process, you can use our SVN repository to get the most up-to-date db4o source code.

Access to the public projects on our Subversion server is available under the following public URL. No login is required.

\href{https://source.db4o.com/db4o/trunk/}{https://source.db4o.com/db4o/trunk/}

The following projects are currently available.

Projects may be under constant development. Source code is not guaranteed to be stable.

Most top-level modules in svn directly map to Eclipse projects, i.e. the root folder contains the Eclipse project metadata.

\begin{itemize}
\item bloat: The byte code manipulation library db4o uses.
\item db4o.cs: The client server implementation for db4o.
\item db4o.cs.optional: Optional client server features for db4o.
\item db4o.instrumentation: The instrumentation basic functionality for db4o.
\item db4o.net: The .NET version of db4o.
\item db4obuild: The build scripts for db4o.
\item db4onqopt: The db4o native query optimizer.
\item db4otaj: The db4o transparent activation/persistence enhancer.
\item db4otools: The db4o tools like Ant scripts.
\item db4ounit: The db4o test project.
\item db4ounit.extensions: Additional unit test functionality.
\item decaf: The Java 1.5 to Java 1.4 converter. You need to checkout all sub-projects on the same level as other db4o projects.
\item drs: The db4o replication system.
\item sharpen: The Java to C\# converter. You need to checkout all sub-projects on the same level as other db4o projects.
\end{itemize}

\subsection{Building Full Distribution}
\label{ref:community-working-with-source-building-full-distribution}

Building full distribution will allow you to get the same db4o packages as you can get from db4o download center. However, the flexibility of the build project also allows you to get only parts of it, like only java distro, only documentation, only tests etc.

The following documentation explains how to build a full distribution using Eclipse version 3.4 Ganymede. It is assumed that you have \href{http://ant.apache.org/}{ant} and one of Eclipse SVN clients (\href{http://subclipse.tigris.org/}{Subclipse} or \href{http://www.eclipse.org/subversive/}{Subversive}) installed.

\subsubsection*{Projects Required}

In order to build db4o you will need to check out the following projects.

\begin{itemize}
\item bloat: The byte code manipulation library db4o uses.
\item db4o.cs: The client server implementation for db4o.
\item db4o.cs.optional: Optional client server features for db4o.
\item db4o.instrumentation: The instrumentation basic functionality for db4o.
\item db4o.net: The .NET version of db4o.
\item db4obuild: The build scripts for db4o.
\item db4onqopt: The db4o native query optimizer.
\item db4otaj: The db4o transparent activation/persistence enhancer.
\item db4otools: The db4o tools like Ant scripts.
\item db4ounit: The db4o test project.
\item db4ounit.extensions: Additional unit test functionality.
\item decaf: The Java 1.5 to Java 1.4 converter. You need to checkout all sub-projects on the same level as other db4o projects.
\item drs: The db4o replication system.
\item sharpen: The Java to C\# converter. You need to checkout all sub-projects on the same level as other db4o projects.
\item doctor: A tool for building the tutorial.
\item tutorial: The db4o tutorial
\end{itemize}

\paragraph*{machine.properties}

You will need to create machine.properties file in db4obuild folder. The contents of the file can be copied from build.xml (see the comments at the beginning of the file). Modify the paths where applicable to set the build variables for your environment.

Read the instructions at top of the build.xml to find out what options are available. Here's an example:

\begin{lstlisting}
file.compiler.jdk1.3=%JAVA_HOME%/bin/javac.exe
file.compiler.jdk1.3.args.optional=-source 1.3
file.jvm.jdk1.5=%JAVA_HOME%/bin/java.exe
dir.workspace=C:/Users/Gamlor/Develop/db4o/db4o-src/
eclipse.home=C:/progs/eclipse
msbuild.executable="C:/Windows/Microsoft.NET/Framework/v4.0.30319/MSBuild.exe" 
\end{lstlisting}

\subsubsection*{Build Preparation}

First you will need to run some preparation scripts. This is done only once per workspace and should not be repeated in the future.

Run build-db4obuild.xml, this will compile some of the tools used in the build process.

You will need to generate a key to sign the tutorial applet. Use the following commands:

keytool -genkey -alias db4objects -keyalg rsa

keytool -export -alias db4objects -file [path]/db4obuild/config/db4objects.crt

Use "kistoa" (without quotes) as your keypass and storepass.

Replace [path] with the path to db4obuild project on your system and make sure that db4objects.crt file is created in db4obuild/config folder.

If you've already generated db4objects key pair before, you will need to delete it before re-generating:

keytool -delete -alias db4objects

You will need to add ant-contrib.jar to your eclipse ant. You can download ant-contib.jar at:

\href{http://sourceforge.net/projects/ant-contrib/}{http://sourceforge.net/projects/ant-contrib}

\begin{itemize}
\item Add ant-contib jar to ant folder in eclipse/plugins.
\item After this is done go to Window-\textgreater{}Preferences menu in Eclipse.
\item Select Ant-\textgreater{}Runtime in the list.
\item Then select "Ant Home Entries".
\item Press "Add External Jar" and select ant-contib.jar location in the plugins folder.
\end{itemize}

\subsubsection*{Running The Build}

Now everything is ready to run db4o build. Right-click build.xml file and select "Run As/Ant Build". You will need to run "buildall" target to generate java and .NET distribution.

\subsection{Patch Submission}
\label{ref:community-working-with-source-patch-submission}

This topic explains how to prepare your patch.

Before writing a patch, please, familiarize yourself with our Coding Style conventions.

You can create a patch using "Create Patch" SVN command.

If you are using \href{http://www.polarion.org/index.php?page=overview&project=subversive}{Subversive} plugin you can use the following steps:

\begin{itemize}
\item select the project in Package Explorer;
\item right-click and select Team/Create Patch;
\item select "Save In File System" and choose the file name;
\item click "Next" and "Finish";
\item check unversioned resources that should be included in the patch;
\item click "OK"
\end{itemize}

Once the patch is created you should register the functionality provided with our \href{http://tracker.db4o.com/secure/Dashboard.jspa}{Jira tracking system} by creating new issue, submitting the description, patch and \hyperref[ref:community-working-with-source-db4o-testing-framework]{test case} (p.~\pageref{ref:community-working-with-source-db4o-testing-framework}) if applicable.

